\subsection{Dimension and flatness}

Let \(\rho: A \to B\) be a ring homomorphism. We denote by (PM) the following condition:

(PM) There exists a B-module  faithfully flat over A such that, for every prime ideal  of B, one has \(_{B}\).

Remarks. --- 1) Condition (PM) is satisfied when there exists a finitely generated B-module, faithfully flat over A and with support equal to Spec(B). This is the case, in particular, if the A-module B is faithfully flat.

\begin{enumerate}
\def\labelenumi{\arabic{enumi})}
\setcounter{enumi}{1}
\item
  The existence of a B-module  faithfully flat over A implies the injectivity of \(\rho\) (I, § 3, no. 5, Proposition 8), and the surjectivity of the map \({}^{a}\rho : \operatorname{Spec}(B) \to \operatorname{Spec}(A)\) (II, § 2, no. 5, Corollary 4 to Proposition 11).
\item
  Suppose that \(\rho: A \to B\) is a local homomorphism of local rings and that there exists a B-module \(N\) flat over A such that  for every prime ideal \(A\)\(N \otimes_B \kappa(\mathfrak{q}) \neq 0\)\(B\) of B. Then \(N\)\(A\) is faithfully flat over A and \(\rho\) therefore enjoys property (PM): indeed, we have , hence \(N/\mathfrak{m}_{B}N = N \otimes_B \kappa(\mathfrak{m}_{B}) \neq 0\) and a fortiori \(N \neq \mathfrak{m}_{B}N\)\(N \neq \mathfrak{m}_{A}N\), and the conclusion follows from Proposition 1 of I, § 3, no. 1.
\end{enumerate}

PROPOSITION 1. --- Let \(\rho: A \to B\) be a ring homomorphism satisfying condition (PM).

\begin{enumerate}
\def\labelenumi{\alph{enumi})}
\item
  Let \(h: \mathbf{A} \to \mathbf{A}'\) be a ring homomorphism. Then the homomorphism \(\rho': \mathbf{A}' \to \mathbf{A}' \otimes_{\mathbf{A}} \mathbf{B}\) deduced from \(\rho\) satisfies condition (PM).
\item
  Let  be a prime ideal of B and \(^{-1}\). The canonical homomorphism \(_{p}\)\(_{q}\) satisfies condition (PM).
\item
  Let  be a prime ideal of B and \(\mathfrak{p} = \rho^{-1}(\mathfrak{q})\). For every prime ideal  of A contained in , there exists a prime ideal \(\mathfrak{p}'\) of B lying over  and contained in \(\mathfrak{q}'\)\(p'\).
\end{enumerate}

Let  be a B-module faithfully flat over A such that  for every prime ideal \(N \otimes_{B} \kappa(q) \neq 0\) of B.

We prove \(a)\). The A'-module \(\mathbf{N}' = \mathbf{A}' \otimes_{\mathbf{A}} \mathbf{N}\) is faithfully flat (I, § 3, no. 3, Proposition 5); let  be a prime ideal of \(\mathbf{B}' = \mathbf{A}' \otimes_{\mathbf{A}} \mathbf{B}\) and  its inverse image in B. We have isomorphisms

\[
\mathrm{N} ^ {\prime} \otimes_ {\mathrm{B} ^ {\prime}} \kappa (\mathfrak {q} ^ {\prime}) \rightarrow \mathrm{N} \otimes_ {\mathrm{B}} \kappa (\mathfrak {q} ^ {\prime}) \rightarrow (\mathrm{N} \otimes_ {\mathrm{B}} \kappa (\mathfrak {q})) \otimes_ {\kappa (\mathfrak {q})} \kappa (\mathfrak {q} ^ {\prime});
\]

since , we also have \(_{B}\)\(_{B'}\).

Let us prove . By Propositions 13 and 14 of II, § 3, no. 4, the -module \(_{p}\) is flat. On the other hand, let \(_{q}\) be a prime ideal of \(_{q}\); it is of the form , where \(_{q}\) is a prime ideal of B contained in \(\mathfrak{b}\) (II, § 3, no. 1, Proposition 3). We have \(\mathfrak{q}\) by the hypothesis on N, and since  is isomorphic to \(_{q}\), we have \(_{q}\)\(_{q}\). Remark 3 allows us to conclude.

Let us prove c). The local homomorphism \(\rho_{\mathfrak{q}}: \mathrm{A}_{\mathfrak{p}} \to \mathrm{B}_{\mathfrak{q}}\) deduced from \(\rho\) satisfies condition (PM) by b). The map \(\operatorname{Spec}(\mathrm{B}_{\mathfrak{q}}) \to \operatorname{Spec}(\mathrm{A}_{\mathfrak{p}})\) is therefore surjective (remark 2), which is what we wanted to prove.

COROLLARY. --- Let F be a closed subset of Spec(A). If Y is an irreducible component of the inverse image of F under the map , then the closure of \(^{a}\rho : \operatorname{Spec}(B) \to \operatorname{Spec}(A)\)\(^{a}\rho(Y)\) is an irreducible component of F.

Let  be an ideal of A such that  and  the prime ideal of B such that \(F = V(a)\). The inverse image under  of F is the subset \(Y = V(\mathfrak{q})\) of , and the closure of  is the irreducible closed subset \(V(\rho(a)B)\) of \(V(\rho^{-1}(q))\).

We need to prove that if  is minimal among the prime ideals of B containing , then \(\rho(a)\) is minimal among the prime ideals of A containing \(\rho^{-1}(q)\). Otherwise, there would exist a prime ideal \(p'\) of A with  and \(a \subset p' \subset \rho^{-1}(q)\); by prop. 1, c), there would exist a prime ideal \(p' \neq \rho^{-1}(q)\) of B such that \(\mathfrak{q}'\) and , whence \(\mathfrak{q}' \subset \mathfrak{q}\) and \(\mathfrak{p}' = \rho^{-1}(\mathfrak{q}')\), contrary to the hypothesis made on \(\rho(a) \subset q' \subset q\)\(q' \neq q\).

PROPOSITION 2. --- Let  be a homomorphism of nonzero rings possessing property (PM). Then we have the inequality

\[
\dim (\mathbf {B}) \geqslant \dim (\mathbf {A}) + \inf _ {\mathfrak {m} \in \mathbf {S}} \dim (\mathbf {B} / \mathfrak {m B})\tag{1}
\]

where S is the set of maximal ideals of A.

We know that \(\dim(A) = \sup_{\mathfrak m \in S} \dim(A_{\mathfrak m})\)\(\S\)\(n^o\) (§ 1, no. 3, Proposition 8). It therefore suffices to establish the inequality

\[
\dim (\mathbf {B}) \geqslant \dim (\mathrm{A} _ {\mathfrak {m}}) + \dim (\mathrm{B} / \mathfrak {m B})\tag{2}
\]

for every maximal ideal m of A. In other words, it is a matter of proving the inequality

\[
\dim (\mathbf {B}) \geqslant n + r\tag{3}
\]

if \(\mathfrak{p}_0\subset \ldots \subset \mathfrak{p}_n\) is a chain of prime ideals of A contained in  and \(\overline{\mathfrak{q}}_0\subset \ldots \subset \overline{\mathfrak{q}}_r\) a chain of prime ideals of . For \(0\leqslant i\leqslant r\), there exists a prime ideal  of B lying over  such that \(\mathfrak{q}_{n+i}\), and \(\overline{\mathfrak{q}}_i = \mathfrak{q}_{n + i} / \mathfrak{mB}\) is a chain of prime ideals of B. Set \(\mathfrak{q}_n\subset \ldots \subset \mathfrak{q}_{n + r}\) for \(\mathfrak{p}_i' = \mathfrak{p}_i\) and \(0\leqslant i\leqslant n-1\), so that \(\mathfrak{p}_n' = \mathfrak m\)\(\mathfrak{p}_0'\subset \ldots \subset \mathfrak{p}_n'\) is a chain of prime ideals of A and \(\mathfrak{q}_n\) lies over \(\mathfrak{p}_n'\). If \(\mathfrak{q}_i\) is a prime ideal of B lying over \(\mathfrak{p}_i'\) (\(1\leqslant i\leqslant n\)), Proposition 1, c), proves that there exists a prime ideal  lying over \(\mathfrak{q}_{i-1}\)\(\mathfrak{p}_{i-1}'\) and contained in \(\mathfrak{q}_i\). By descending induction, one therefore constructs a chain \(\mathfrak{q}_0\subset \ldots \subset \mathfrak{q}_n\) of prime ideals of B such that \(\mathfrak{q}_i\) lies over \(\mathfrak{p}_i'\) for \(0\leqslant i\leqslant n\). Since \(\mathfrak{q}_0\subset \ldots \subset \mathfrak{q}_{n+r}\) is a chain of prime ideals of B, we have proved inequality (3).

Remark 4. --- Let \(\rho: A \to B\) be a local homomorphism of Noetherian local rings satisfying condition (PM). We shall see later (§ 3, no. 4, prop. 7) that in this case equality holds in (1). In the general case there may be strict inequality (p. 84, exercise 1).

COROLLARY. --- For every ideal  of A, one has 
 be a prime ideal of B containing 
, and \(\rho(a)\mathbf{B}\). By prop. 1, the local homomorphism \(p = \rho^{-1}(q)\)\(\rho_{\mathfrak{q}}: A_{\mathfrak{p}} \to B_{\mathfrak{q}}\) deduced from \(\rho\) satisfies (PM), and we therefore have  by prop. 2. By prop. 7 of § 1, no. 3, we have  since \(\dim(A_{\mathfrak{p}}) \leqslant \dim(B_{\mathfrak{q}})\) contains , whence \(\leqslant \dim(A_p)\) for every prime ideal  of B containing \(\leqslant \dim(B_q)\)\(\rho(\mathfrak{a})\mathbf{B}\). The corollary then follows from loc. cit.

Lemma 1. — Let \(\rho: \mathbf{A} \to \mathbf{B}\) be a ring homomorphism and \(\mathfrak{p}\) a prime ideal of A. The continuous map \(^{a}h: \operatorname{Spec}(\mathbf{B} \otimes_{\mathbf{A}} \kappa(\mathfrak{p})) \to \operatorname{Spec}(\mathbf{B})\), associated with the canonical homomorphism \(h:\mathbf{B}\to\mathbf{B}\otimes_{\mathbf{A}}\kappa(\mathfrak{p})\), induces a homeomorphism from \(\operatorname{Spec}(\mathbf{B}\otimes_{\mathbf{A}}\kappa(\mathfrak{p}))\) onto the subspace \((^{a}\mathfrak{p})^{-1}(\mathfrak{p})\) of \(\operatorname{Spec}(\mathbf{B})\) consisting of the prime ideals of \(\mathbf{B}\) above \(\mathfrak{p}\).

The homomorphism \(h\) is composed of the quotient homomorphism from B to \(\mathbf{B} / \rho(\mathfrak{p})\) B and of the canonical homomorphism from \(\mathbf{B} / \rho(\mathfrak{p})\) B into its ring of fractions \((\rho(\mathbf{A} - \mathfrak{p}))^{-1}(\mathbf{B} / \rho(\mathfrak{p}) \mathbf{B})\) . By the remark and the corollary to prop. 13 of II, § 4, no. 3, \(^a h\) therefore induces a homeomorphism of \(\text{Spec}(\mathbf{B} \otimes_{\mathbf{A}} \kappa(\mathfrak{p}))\) onto the subspace of \(\text{Spec}(\mathbf{B})\) formed by the prime ideals \(\mathfrak{q}\) of B which contain \(\rho(\mathfrak{p})\) and are disjoint from \(\rho(\mathbf{A} - \mathfrak{p})\) , that is, which lie over \(\mathfrak{p}\) .

Remark 5. --- By prop. 2 and lemma 1, we therefore have, under the hypotheses of prop. 2, the inequality

\[
\dim (\operatorname{Spec} (B)) \geqslant \dim (\operatorname{Spec} (A)) + \inf _ {\mathfrak {p} \in \operatorname{Spec} (A)} \dim \left(({}^{a}\rho)^{-1} (\mathfrak {p})\right).\tag{4}
\]

